\chapter*{Summary}
\addcontentsline{toc}{chapter}{Summary}

Imagine standing at the seashore: each wave arrives, retreats, and is replaced by another, yet the ocean itself has no beginning or end. Human perception is drawn to the rhythm, to the measurable intervals of rise and fall, but the ocean's essence is not bound by these temporal markers. In a similar way, this work proposes a model of Artificial Intelligence that seeks to transcend the anthropocentric constraints of sequence, space, and time.

We introduce the Apeiron Framework\footnote{Named after the ancient Greek philosophical concept of \textit{to apeiron} ($\tau\grave{o}$ $\ddot{\alpha}\pi\varepsilon\iota\rho o\nu$), introduced by Anaximander (c. 610--546 BCE) to denote the boundless, indefinite, and eternal primordial substance from which all things arise and to which they return.}, a seventeen-layer theoretical architecture for a fundamentally different kind of AI---one designed not to replicate human knowledge, but to autonomously construct its own body of knowledge from elementary observables. Rather than a rigid taxonomy, the seventeen layers constitute a dynamic field of relational densities, where entities can exist across multiple layers, outside linear temporality, and in parallel frames of interpretation.

The framework begins at Layer~1 with a field of irreducible, context-independent observables---entities defined without human semantic labeling. Through successive layers of relational emergence, functional organization, adaptive meta-learning, and ontological creation, the system rebuilds categories, temporalities, and causal structures from first principles. The human role shifts from instructor to observer: we do not tell the system what to see; we provide the initial substrate and then witness what knowledge it autonomously generates, offering humanity a new perspective on the latent structure of the domains it observes.

We provide a rigorous mathematical specification for the foundational Layer~1 (the \textit{Epistemic Singularity}), including decomposition operators grounded in Boolean algebra, measure theory, category theory, and information theory. A reference implementation in Python validates the framework through a falsifiable experiment, confirming the multi-axial convergence of atomicity frameworks on a fundamental irreducible unit. The seventeen layers thus function both as scaffolding for thought and as apertures into the infinite---like windows in a cathedral that illuminate, but never contain, the light that passes through them.

This extended version provides complete formal proofs of all key propositions, a comprehensive analysis of computational complexity, a detailed ethical and governance framework, and a suite of falsifiable predictions that chart the course for empirical validation of the entire architecture.